\documentclass[10pt,aspectratio=169]{beamer}
% Preámbulo modular para presentaciones Beamer (Modelación Estadística)
% Diseño y paleta institucional del Tecnológico de Monterrey

\documentclass[aspectratio=169,xcolor={dvipsnames,table}]{beamer}

\usepackage[utf8]{inputenc}
\usepackage[spanish,mexico]{babel}
\usepackage{amsmath,amssymb,amsfonts,mathtools}
\usepackage{graphicx}
\usepackage{listings}
\usepackage{booktabs}
\usepackage{tikz}
\usetikzlibrary{matrix,arrows,decorations.pathmorphing,shapes.geometric,positioning}

% Paleta Institucional Tecnológico de Monterrey
\definecolor{TecAzul}{HTML}{0039A6}       % Azul TEC: color primario institucional
\definecolor{TecAzulOscuro}{HTML}{1A2E51} % Azul marino oscuro
\definecolor{TecRojo}{HTML}{EC2661}       % Rojo de acento (paleta miTec)
\definecolor{TecGris}{HTML}{646464}       % Gris neutro para comentarios y subtítulos
\definecolor{TecFondoBloque}{HTML}{F4F6F9}% Fondo suave para bloques

% Configuración de Tema Beamer
\usetheme{Boadilla}
\usecolortheme[named=TecAzulOscuro]{structure}
\setbeamercolor{alerted text}{fg=TecRojo}
\setbeamercolor{palette primary}{bg=TecAzulOscuro,fg=white}
\setbeamercolor{palette secondary}{bg=TecAzul,fg=white}
\setbeamercolor{palette tertiary}{bg=TecAzulOscuro!80!black,fg=white}
\setbeamercolor{title}{fg=TecAzulOscuro,bg=TecFondoBloque}
\setbeamercolor{frametitle}{fg=TecAzulOscuro,bg=TecFondoBloque}

% Bloques personalizados elegantes
\setbeamercolor{block title}{bg=TecAzulOscuro,fg=white}
\setbeamercolor{block body}{bg=TecFondoBloque,fg=black}
\setbeamercolor{block title alerted}{bg=TecRojo,fg=white}
\setbeamercolor{block body alerted}{bg=TecRojo!8,fg=black}
\setbeamercolor{block title example}{bg=TecAzul,fg=white}
\setbeamercolor{block body example}{bg=TecAzul!8,fg=black}

% Redondear bordes y añadir sombra sutil
\setbeamertemplate{blocks}[rounded][shadow=true]
\setbeamertemplate{navigation symbols}{}

% Pie de página limpio con número de diapositiva
\setbeamertemplate{footline}{
  \leavevmode%
  \hbox{%
  \begin{beamercolorbox}[wd=.7\paperwidth,ht=2.25ex,dp=1ex,left]{author in head/foot}%
    \hspace*{2ex}\insertshortauthor~~(\insertshortinstitute) \hfill \insertshorttitle
  \end{beamercolorbox}%
  \begin{beamercolorbox}[wd=.3\paperwidth,ht=2.25ex,dp=1ex,right]{date in head/foot}%
    \insertshortdate{}\hspace*{2em}
    \insertframenumber{} / \inserttotalframenumber\hspace*{2ex}
  \end{beamercolorbox}}%
  \vskip0pt%
}

% Configuración de Listings de Python para visualización pedagógica de código
\lstset{
    language=Python,
    basicstyle=\ttfamily\footnotesize,
    keywordstyle=\color{TecAzulOscuro}\bfseries,
    stringstyle=\color{TecRojo},
    commentstyle=\color{TecGris}\itshape,
    showstringspaces=false,
    breaklines=true,
    frame=lines,
    rulecolor=\color{TecAzulOscuro!30},
    backgroundcolor=\color{TecFondoBloque},
    aboveskip=0.5em,
    belowskip=0.5em,
    literate={á}{{\'a}}1 {é}{{\'e}}1 {í}{{\'i}}1 {ó}{{\'o}}1 {ú}{{\'u}}1
             {Á}{{\'A}}1 {É}{{\'E}}1 {Í}{{\'I}}1 {Ó}{{\'O}}1 {Ú}{{\'U}}1
             {ñ}{{\~n}}1 {Ñ}{{\~N}}1 {¿}{{?`}}1 {¡}{{!`}}1
}

% Metadatos institucionales por defecto
\title[Modelación Estadística]{Modelación Estadística para Ciencia de Datos e Ingeniería}
\author[J. Castillo Colmenares]{Juliho David Castillo Colmenares}
\institute[Tec de Monterrey]{Tecnológico de Monterrey \\ Escuela de Ingeniería y Ciencias}
\date{\today}
% Comandos y macros estadísticas compartidas para las presentaciones Beamer (Metropolis)

% Conjuntos numéricos básicos
\newcommand{\R}{\mathbb{R}}
\newcommand{\N}{\mathbb{N}}
\newcommand{\Q}{\mathbb{Q}}
\newcommand{\Z}{\mathbb{Z}}
\newcommand{\set}[1]{\left\{ #1 \right\}}
\newcommand{\abs}[1]{\left|#1\right|}
\newcommand{\norm}[1]{\left\| #1 \right\|}

% Comandos estadísticos y probabilísticos del curso
\newcommand{\Var}{\operatorname{Var}}
\newcommand{\cov}{\operatorname{Cov}}
\newcommand{\corr}{\rho}
\newcommand{\comb}[2]{\begin{pmatrix} #1 \\ #2 \end{pmatrix}}
\newcommand{\s}{\sigma}
\newcommand{\particion}{\mathfrak{P}}
\newcommand{\card}[1]{\#\left( #1 \right)}
\newcommand{\seq}[3]{\set{#1}_{#2}^{#3}}
\newcommand{\E}{\mathbb{E}}
\newcommand{\Prob}{\mathbb{P}}

% Entornos pedagógicos y cajas en español académico natural
\newenvironment{discusion}{%
  \begin{alertblock}{Pregunta de Discusión}%
}{%
  \end{alertblock}%
}

\newenvironment{reflexion}{%
  \begin{alertblock}{Reflexión para el Aula}%
}{%
  \end{alertblock}%
}

\newenvironment{paradoja}{%
  \begin{alertblock}{Paradoja de Exploración}%
}{%
  \end{alertblock}%
}

\newenvironment{motivacion}{%
  \begin{exampleblock}{Motivación: Ciencia de Datos en la Práctica}%
}{%
  \end{exampleblock}%
}

\newenvironment{retoclase}{%
  \begin{block}{Reto Rápido en Equipo}%
}{%
  \end{block}%
}

\title[Varias proporciones]{Sección 07.12: Pruebas de varias proporciones}\subtitle{Contraste global y comparaciones posteriores}
\author[J. Castillo Colmenares]{Juliho Castillo Colmenares}\institute{Tecnológico de Monterrey}\date{}
\begin{document}
\begin{frame}[plain]\titlepage\end{frame}
\begin{frame}{Hoja de ruta}\small\begin{enumerate}\item Contraste global.\item Estadística combinada.\item Procedimiento de Marascuilo.\end{enumerate}\end{frame}
\begin{frame}{Fuente y alcance}\small Esta sección trata $k$ poblaciones independientes con una respuesta binaria: éxito o fracaso.\begin{block}{Pregunta guía}¿Son iguales todas las proporciones de éxito?\end{block}\end{frame}
\begin{frame}{Hipótesis global}\small\[H_0:p_1=p_2=\cdots=p_k\qquad\text{frente a}\qquad H_a:\text{al menos dos difieren}.\]\end{frame}
\begin{frame}{Datos}\small En la muestra $j$ se observan $X_j$ éxitos de $n_j$ ensayos y $\hat p_j=X_j/n_j$.\end{frame}
\begin{frame}{Proporción global}\small Bajo $H_0$, la proporción común se estima con\[\bar p=\frac{\sum_jX_j}{\sum_jn_j}.\]\end{frame}
\begin{frame}{Estadística global}\small\[\chi^2=\sum_{j=1}^{k}\frac{(X_j-n_j\bar p)^2}{n_j\bar p(1-\bar p)}.\]\end{frame}
\begin{frame}{Grados de libertad}\small El contraste global se compara con $\chi^2_{k-1}$. Un rechazo solo indica que no todas las proporciones son iguales.\end{frame}
\begin{frame}{Por qué se requieren comparaciones}\small La prueba global no identifica los pares que difieren. Comparar todos los pares sin ajuste inflaría el error tipo I.\end{frame}
\begin{frame}{Procedimiento de Marascuilo}\small El método compara cada diferencia $|\hat p_i-\hat p_j|$ con un rango crítico $r_{ij}$ después de un rechazo global.\end{frame}
\begin{frame}{Rango crítico}\small\[r_{ij}=\sqrt{\chi^2_{\alpha,k-1}}\sqrt{\frac{\hat p_i(1-\hat p_i)}{n_i}+\frac{\hat p_j(1-\hat p_j)}{n_j}}.\]\end{frame}
\begin{frame}{Procedimiento I}\small\begin{enumerate}\item Calcular cada $\hat p_j$.\item Calcular $\bar p$.\item Evaluar $\chi^2$ global.\end{enumerate}\end{frame}
\begin{frame}{Procedimiento II}\small\begin{enumerate}\setcounter{enumi}{3}\item Si el global no rechaza, detenerse.\item Si rechaza, calcular rangos $r_{ij}$.\item Reportar pares con diferencia mayor al rango.\end{enumerate}\end{frame}
\begin{frame}{Condiciones}\small Las muestras deben ser independientes y las aproximaciones binomiales razonables. El diseño fija $n_j$ por población.\end{frame}
\begin{frame}{Aplicación: datos}\small Para cada grupo registre $X_j$, $n_j$ y $\hat p_j$. Ordenar los grupos ayuda a leer las diferencias, pero no cambia el contraste.\end{frame}
\begin{frame}{Aplicación: contraste global}\small Sustituya los conteos en la estadística y compare con $\chi^2_{k-1}$ o calcule el valor $p$.\end{frame}
\begin{frame}{Aplicación: decisión global}\small Un valor $p$ pequeño indica heterogeneidad entre las proporciones, sin señalar todavía los pares responsables.\end{frame}
\begin{frame}{Aplicación: pares}\small Para cada $i<j$, calcule la diferencia absoluta y el rango crítico de Marascuilo.\end{frame}
\begin{frame}{Errores frecuentes}\small\begin{itemize}\item Hacer comparaciones por pares sin contraste global.\item Omitir el ajuste de Marascuilo.\item Confundir heterogeneidad global con todos los pares distintos.\end{itemize}\end{frame}
\begin{frame}{Plantilla de reporte}\small\begin{block}{Reporte} ``El contraste global dio $\chi^2=\cdots$, $gl=k-1$, $p=\cdots$. Tras el ajuste, los pares con diferencia detectable son [lista].''\end{block}\end{frame}
\begin{frame}{Síntesis}\small\begin{itemize}\item La proporción global define el contraste bajo igualdad.\item $\chi^2$ responde la pregunta global.\item Marascuilo localiza diferencias sin inflar el error.\end{itemize}\end{frame}
\begin{frame}{Conexión con el capítulo}\small Esta sección extiende las pruebas de una y dos proporciones al escenario de $k$ grupos independientes.\end{frame}
\end{document}
